同一份数据表，同一个回归模型，两个人可以得出完全相容的两句话：``用了这个功能的用户留存率高 \(12\) 个百分点''，以及``把这个功能推给所有人，留存率不会变''。前一句描述的是看见什么，后一句描述的是做了什么之后会怎样。它们都可能为真，因为它们回答的根本不是同一个问题。

预测问的是：观察到某些特征之后，结果通常怎样分布。因果问的是：主动改变一个条件之后，结果会怎样分布。第一个问题的答案由联合分布决定，数据够多就能逼近；第二个问题的答案不在联合分布里——同一个联合分布可以由多套机制生成，它们对干预的回答彼此矛盾，而任何检验都区分不了它们。差额必须由结构假设补上，这些假设写在图上、写在设计里，就是写不进数据表。

本章不把因果推断当作一类新算法，而是把它拆成三层：先定义想回答的因果量，再判断它能否由现有数据识别，最后才挑估计方法。这个顺序不能颠倒。若识别所需的无混杂、正值性或一致性不成立，加样本、换更灵活的模型、把标准误算得更细，都补不上目标与数据之间那道裂缝——它们只会让一个错误的量被估得更精确。

六节沿这条链依次推进。第一节把观察、干预与反事实分成三个层级，说明为什么要先固定估计目标再谈计算。第二节引入结构因果模型，把干预写成机制替换，得到截断分解——干预删掉一个因子，条件化只是重新归一化，两种操作的算术差别就在这里。第三节回答该控制谁：控制混杂消除偏差，控制碰撞点制造偏差，控制中介截断效应，后门准则给出判据。第四节换用潜在结果的语言，把效应定义在同一个体的两个世界之间，逐条交代一致性、可忽略性、正值性各挡住什么。第五节进入估计，把结果回归、倾向加权与增广估计放进同一个公式，说明双重稳健给的是两次机会而非豁免权。第六节收束边界：识别出一个量，既不表示它在别的人群成立，也不表示按它行动是对的。

各种方法要付出什么，值得先摆明。随机化最强，因为分配机制已知，但很多问题不能随机——不能指派谁去吸烟，不能指派谁失业。工具变量绕开了这一点，要认下的是排除性约束，那是一句关于不存在某条路径的断言，数据检验不了，而且识别出来的常常只是依从者子群上的局部效应。双重稳健允许两个辅助模型错一个，两个都错就没有第三条命。没有哪种方法能把观察数据变成实验，它们只是在不同的地方付账。

本章主要讨论二元处理与平均处理效应，因为这个最简单的设定已经足以暴露全部结构性困难。连续处理、纵向干预、断点回归、因果发现都能在同一框架里展开，但每种设计都要自己那份识别论证，不是把二元公式换个符号就成立。

\subsection{怎么读这一章}

顺读的话，先看每节开头那个具体场景，再沿``问题—识别—估计—迁移—决策''往下走：第一节确定要问什么，第二至第四节判断观察数据能不能回答，第五节才讨论怎么估，第六节检查结论能不能用在目标环境里。想弄清公式从哪来，重点在机制替换、后门调整、潜在结果与增广估计这四处。

如果手上正有一份分析要做，倒着读更省时间：从第六节的报告结构开始逐项自查，任何一项说不清楚，都不该指望换一个更复杂的估计器把它糊过去。

\subsection{本章篇目}

以下篇目按概念依赖排列；若从具体困惑进入，也可以直接打开最接近的一篇，再沿篇内链接回补。

\begin{itemize}
\tightlist
\item \hyperref[sec:13-Machine-Learning/12-Causal-Inference/01]{``观察、干预与反事实回答不同问题''}；
\item \hyperref[sec:13-Machine-Learning/12-Causal-Inference/02]{``结构因果模型把干预写成机制替换''}；
\item \hyperref[sec:13-Machine-Learning/12-Causal-Inference/03]{``混杂、碰撞点与后门调整''}；
\item \hyperref[sec:13-Machine-Learning/12-Causal-Inference/04]{``潜在结果把处理效应定义在替代世界之间''}；
\item \hyperref[sec:13-Machine-Learning/12-Causal-Inference/05]{``回归、加权与双重稳健估计各自依赖什么''}；
\item \hyperref[sec:13-Machine-Learning/12-Causal-Inference/06]{``可识别不等于可外推也不等于可决策''}。
\end{itemize}

图与观察数据之间的界线——两张图什么时候蕴涵完全相同的条件独立关系、观察数据最多定到哪一层——由\hyperref[sec:12-Real-Analysis-Support/10-Essential-Structures/04]{``因果结构的代数形式''}一节交代；本章从那条界线往下走，默认图或设计已经给定，追问的是在此之上还能算出什么、又该在哪里停手。
